% Part: first-order-logic
% Chapter: model-theory
% Section: reducts-and-expansions

\documentclass[../../../include/open-logic-section]{subfiles}

\begin{document}

\olfileid{mod}{bas}{red}
\section{Symbolen weglaten of toevoegen (reducten en expansies)}

Het is vaak handig om talen met gedeelde symbolen te vergelijken, en
ook de !!{structure}s voor die talen. Meestal gaat het om deze
situatie: alle symbolen uit !!a{language}~$\Lang{L}$ komen ook voor in
!!a{language}~$\Lang{L'}$, dus $\Lang{L} \subseteq \Lang{L'}$. Van een
$\Lang{L}$-!!{structure}~$\Struct{M}$ kun je dan altijd een
$\Lang{L'}$-!!{structure} maken. Je geeft de extra symbolen een
interpretatie, dus een betekenis binnen de structuur, en laat de
interpretaties van de gedeelde symbolen hetzelfde. Andersom kun je van
een $\Lang{L'}$-structuur~$\Struct{M'}$ weer een
$\Lang{L}$-structuur maken. Je kijkt dan niet meer naar de
interpretaties van de symbolen die niet in~$\Lang{L}$ voorkomen.

\begin{defn}
\ollabel{defn:reduct}
Stel dat $\Lang L \subseteq \Lang L'$. Laat $\Struct M$ een
$\Lang L$-!!{structure} zijn en $\Struct M'$ een
$\Lang L'$-!!{structure}. We noemen $\Struct M$ het
\emph{reduct} van $\Struct M'$ tot $\Lang L$. Omgekeerd noemen we
$\Struct M'$ een \emph{expansie} van $\Struct M$ tot $\Lang L'$.
Dat bedoelen we precies wanneer aan de volgende vier voorwaarden is
voldaan:
\begin{enumerate}
\item De twee structuren hebben hetzelfde domein:
  $\Domain{M} = \Domain{M'}$
\item Elke !!{constant}~$c \in \Lang L$ krijgt in beide structuren
  dezelfde interpretatie: $\Assign{c}{M} =
  \Assign{c}{M'}$.
\item Elke !!{function}~$f \in \Lang L$ krijgt in beide structuren
  dezelfde interpretatie: $\Assign{f}{M} =
  \Assign{f}{M'}$.
\item Elk !!{predicate}~$P \in \Lang L$ krijgt in beide structuren
  dezelfde interpretatie: $\Assign{P}{M} =
  \Assign{P}{M'}$.
\end{enumerate}
\end{defn}

\begin{prop}
\ollabel{prop:reduct}
Als een $\Lang{L}$-!!{structure}~$\Struct{M}$ een reduct is van een
$\Lang{L'}$-!!{structure}~$\Struct{M'}$, dan heeft elke
$\Lang{L}$-!!{sentence}~$!A$ in beide structuren dezelfde
waarheidswaarde:
\[
\Sat{M}{!A} \text{ iff } \Sat{M'}{!A}.
\]
\end{prop}

\begin{proof}
  Oefening.
\end{proof}

\begin{prob}
Bewijs de propositie \olref[mod][bas][red]{prop:reduct}.
\end{prob}

\begin{defn}
Neem een $\Lang{L}$-structuur $\Struct{M}$. De grotere taal
$\Lang{L'} = \Lang{L} \cup \{P\}$ is de expansie van $\Lang{L}$ die
je maakt door precies één $n$-plaatsig !!{predicate}~$P$ toe te voegen.
Laat
$R \subseteq \Domain{M}^n$ een $n$-plaatsige relatie zijn. Dan
schrijven we $\Expan{M}{R}$ voor de expansie~$\Struct{M'}$
van~$\Struct{M}$. In die expansie interpreteren we het nieuwe predicaat
als die relatie, dus $\Assign{P}{M'} = R$.
\end{defn}

\end{document}
